\documentclass[a4paper,10pt,fontset=none]{ctexart}
\usepackage{amsmath,amssymb}
\usepackage{hyperref}
\hypersetup{pdftitle={一线去化疗证据强化，关注套细胞淋巴瘤新方案审批与放量},pdfauthor={行业研究},bookmarksnumbered=true}
% Reusable preamble/header for Chinese research PDFs.
% Compile with XeLaTeX and keep ruc-logo.pdf beside the document.
\usepackage[a4paper,top=25mm,left=25mm,bottom=25mm,right=25mm,headheight=35pt,headsep=5mm,footskip=10mm]{geometry}
\usepackage{fontspec}
\usepackage{xeCJK}
\usepackage{graphicx}
\usepackage{fancyhdr}
\usepackage{eso-pic}
\usepackage{booktabs}
\usepackage{longtable}
\usepackage{array}
\usepackage{tabularx}
\usepackage{etoolbox}
\usepackage{enumitem}
\usepackage{ragged2e}
\usepackage{xcolor}
\usepackage{microtype}
\usepackage{needspace}
\usepackage{indentfirst}
\setmainfont{Times New Roman}
\setsansfont{Arial}
\setmonofont{Courier New}
\setCJKmainfont[AutoFakeBold=1.5]{Songti SC}
\setCJKsansfont[AutoFakeBold=1.5]{Songti SC}
\setCJKmonofont{Songti SC}
\newCJKfontfamily\songti[AutoFakeBold=1.5]{Songti SC}
\newCJKfontfamily\heiti[AutoFakeBold=2]{Songti SC}
\newCJKfontfamily\fangsong{Songti SC}
\setlength{\parindent}{2em}
\setlength{\parskip}{0pt}
\linespread{1.5}
\AtBeginDocument{\zihao{-4}\songti}
\setlength{\emergencystretch}{3em}
% Compact headings adapted from the user's reference paper; manual numbering remains in the titles.
\ctexset{
  section={format=\centering\songti\zihao{-3},beforeskip=1.0ex,afterskip=0.5ex},
  subsection={format=\raggedright\songti\zihao{-4},beforeskip=1.0ex,afterskip=0.5ex},
  subsubsection={format=\raggedright\songti\zihao{-4},beforeskip=1.0ex,afterskip=0.5ex}
}
\setlist{itemsep=0pt,topsep=0pt,parsep=0pt,partopsep=0pt}
\setlist[enumerate,1]{label=\arabic*.,leftmargin=2em}
\pagestyle{fancy}
\fancyhf{}
\fancyhead[C]{\includegraphics[pagebox=mediabox,width=5.2cm]{ruc-logo.pdf}}
\fancyfoot[C]{\zihao{-5}\thepage}
\renewcommand{\headrulewidth}{0.4pt}
\fancypagestyle{plain}{\fancyhf{}\fancyhead[C]{\includegraphics[pagebox=mediabox,width=5.2cm]{ruc-logo.pdf}}\fancyfoot[C]{\zihao{-5}\thepage}\renewcommand{\headrulewidth}{0.4pt}}
\fancypagestyle{empty}{\fancyhf{}\fancyhead[C]{\includegraphics[pagebox=mediabox,width=5.2cm]{ruc-logo.pdf}}\fancyfoot[C]{\zihao{-5}\thepage}\renewcommand{\headrulewidth}{0.4pt}}
\setlength{\tabcolsep}{5pt}
\renewcommand{\arraystretch}{1.0}
\AtBeginEnvironment{longtable}{\zihao{5}\linespread{1.0}\selectfont\setlength{\LTpre}{0.5\baselineskip}\setlength{\LTpost}{0.5\baselineskip}\renewcommand{\arraystretch}{1.0}}
\definecolor{linknavy}{RGB}{36,70,105}
\hypersetup{colorlinks=true,linkcolor=linknavy,urlcolor=linknavy,citecolor=linknavy}
\newcommand{\RUCTableCaption}[1]{\Needspace{3\baselineskip}\par\vspace{0.5\baselineskip}\begin{center}{\heiti\zihao{5}\textbf{#1}\par}\end{center}\vspace{0.5\baselineskip}}
\newcommand{\RUCTableAfterSpace}{\par\vspace{0.7\baselineskip}}
\newcommand{\RUCNote}[1]{\par\vspace{0.35\baselineskip}\noindent{\zihao{5}\songti\textbf{#1}\par}\vspace{0.7\baselineskip}}
% 参考文献条目正文与文章正文统一为小四号宋体；行距沿用全局1.5倍设置。
% 首条参考文献前增加半个正文行距，使标题与正文间距约为原来的两倍；
% 后续条目不增加额外间距。
\newif\ifRUCFirstReference
\RUCFirstReferencetrue
\newcommand{\RUCReference}[2]{\Needspace{8\baselineskip}\ifRUCFirstReference\vspace{0.5\baselineskip}\RUCFirstReferencefalse\fi\par\noindent\hangindent=2em\hangafter=1{{\zihao{-4}\songti #1\ #2}\par}}
\newcommand{\RUCSubtitle}{}
\newcommand{\RUCDateLine}{}
\makeatletter
\renewcommand{\maketitle}{%
  \par\begingroup\centering
  {\fontsize{22}{26.4}\selectfont\songti\textbf{\@title}\par}
  \ifx\RUCSubtitle\empty\else\vspace{10pt}{\fontsize{15}{18}\selectfont\songti\RUCSubtitle\par}\fi
  \ifx\RUCDateLine\empty\else\vspace{8pt}{\zihao{5}\songti\RUCDateLine\par}\fi
  \endgroup
  \thispagestyle{plain}
  \vspace{2pt}
}
\makeatother

\pagestyle{fancy}
\title{一线去化疗证据强化，关注套细胞淋巴瘤新方案审批与放量}
\renewcommand{\RUCSubtitle}{2026年ASH会议前瞻}
\renewcommand{\RUCDateLine}{医药生物行业专题研究｜2026年11月10日}
\author{}
\date{}
\begin{document}
\maketitle

\section*{投资要点}
\noindent\textbf{核心判断：}MANGROVE三期研究的摘要主要分析显示，泽布替尼联合利妥昔单抗（ZR）较苯达莫司汀联合利妥昔单抗（BR）显著改善既往未治疗套细胞淋巴瘤（MCL）患者的无进展生存期（PFS），为一线去化疗方案提供随机对照证据；美国FDA已授予补充新药申请优先审评，欧盟与中国亦已进入申报审评流程。若获批，商业兑现取决于可治疗人群边界、持续治疗时间、实际净价和支付准入。优先审评仅改变审评时限，不代表审评结论或获批保证。\par
\noindent\textbf{证据进展：}ASH官网摘要披露510例入组、中位随访44.1个月，独立评审PFS风险比（HR）0.57；ZR组中位PFS尚未达到，对照组39.2个月。4年PFS率分别为62.4\%和46.2\%。OS在本次分析时仍不成熟，虽有数值趋势，尚不足以确认生存获益。\par
\noindent\textbf{投资观察：}泽布替尼已具备大体量全球销售基数。公司披露2026年第三季度BRUKINSA全球销售额约14亿美元，同比增长34\%；因此，一线MCL属于潜在新增适应症，不能与既有适应症销售混算。关注12月13日ASH全会正式展示是否更新疗效、安全性、患者报告结局及亚组数据，以及FDA、EMA、NMPA后续审评节点。\par
\noindent\textbf{主要待验证条件：}全会及后续随访确认获益稳健性；审评机构接受临床获益与风险证据；真实世界患者适用范围与临床路径清晰；长期BTK抑制剂暴露的耐受性和停药管理可接受；各市场医保、商业保险与定价支持放量。

\section*{一、行业事件与治疗需求}
\subsection*{（一）一线治疗的核心人群与未满足需求}
MCL是成熟B细胞非霍奇金淋巴瘤亚型，患者以老年人为主，疾病生物学和临床过程存在较大异质性。低肿瘤负荷、无症状患者可在专科评估后观察；需要系统治疗者则按年龄、体能、移植适合度、TP53状态、增殖指数和疾病风险分层。移植适合的年轻患者常采用强化免疫化疗及巩固策略；年长或不适合自体造血干细胞移植（ASCT）的患者，常见一线方案包括BR等免疫化疗。治疗选择并非单一路径，不能把全部新诊断患者视为同一可及市场。

传统免疫化疗可带来持续缓解，但烷化剂、细胞毒药物相关的骨髓抑制、感染及输注负担，对高龄或合并症患者尤其重要。BTK抑制剂改变了复发难治MCL治疗格局；一线联合治疗进一步探索在保留疾病控制的同时减少化疗暴露。MANGROVE的ZR方案前6个周期联合利妥昔单抗，之后继续口服泽布替尼至疾病进展或不能耐受；这不是固定疗程方案，意味着潜在治疗便利与长期暴露、累积成本和持续安全性需共同衡量。

\subsection*{（二）临床证据与全球销售之间的传导}
该事件的商业逻辑可拆为四个连续环节：随机对照研究证实临床获益；监管机构认可适应症和用药标签；临床指南、医生与患者接受新方案；支付体系覆盖并形成持续用药。任何一环均可能延迟或压缩收入兑现。MCL为相对罕见的血液肿瘤，单一适应症对公司收入影响需放在泽布替尼既有大基数中衡量，且不能将风险比直接换算为市场份额或销售增速。

\section*{二、MANGROVE数据解读}
\subsection*{（一）研究设计及入组边界}
MANGROVE（BGB-3111-306；NCT04002297）是一项全球、多中心、开放标签、随机三期研究。510例既往未治疗MCL患者按1:1分配至ZR组（257例）和BR组（253例）。入组年龄为70岁及以上，或60至未满70岁且不适合ASCT；ECOG体能状态为0至2，并须有可测量病灶。ZR组接受泽布替尼联合利妥昔单抗6个周期，随后泽布替尼单药直至进展；BR组接受6个周期后观察。两组利妥昔单抗均为每月给药，共6次，未安排维持治疗。主要终点为独立评审委员会（IRC）评估PFS。

该设计主要代表年长或不适合移植患者，不能直接外推至所有年轻、移植适合人群。开放标签可能影响部分评估行为，但主要终点由IRC评估。研究对照为BR，而非当前所有地区和风险分层下的全部标准方案；治疗路径需结合本地指南及患者个体特征判断。

\subsection*{（二）疗效信号强，生存结论仍待成熟}
ASH摘要报告中位随访44.1个月。IRC评估的PFS HR为0.57（95\%置信区间0.43至0.76，单侧P<0.0001），ZR组中位PFS尚未达到（95\%置信区间下限56.5个月，上限未估计），BR组为39.2个月（95\%置信区间33.5至55.7个月）；4年PFS率为62.4\%和46.2\%。预设亚组包括高危MIPI、母细胞样/多形性组织学、Ki-67不低于30\%及TP53突变，摘要称获益方向一致。亚组样本较小，仍需正式报告中的交互作用检验、置信区间和缺失数据说明。

OS HR为0.71（95\%置信区间0.51至0.99，单侧P=0.0222），4年OS率为75.1\%和67.6\%，但OS仍被研究者界定为不成熟。对照组后续接受抗癌治疗者中88.2\%使用过BTK抑制剂，这可能影响组间OS差异解释。当前可支持“PFS改善且观察到OS数值趋势”，不能写成“已证实延长总生存”。

\subsection*{（三）安全性需按时间窗口理解}
安全性结果取决于观察时间和治疗暴露。摘要报告，从开始治疗至疾病进展或开始新抗癌治疗的安全性报告期内，ZR与BR组≥3级不良事件发生率分别为75.0\%和73.6\%，严重不良事件分别为56.6\%和46.0\%，导致死亡的不良事件分别为9.0\%和8.8\%。而仅观察前6个月诱导期时，≥3级不良事件为39.8\%和58.4\%，严重不良事件为19.1\%和23.6\%。两种窗口回答的问题不同；鉴于ZR组持续用药、BR组完成6周期后观察，累计发生率不宜被解读为疗程等长的直接毒性比较。

按每100人月暴露校正后，摘要列示感染、嗜中性粒细胞减少、贫血、血小板减少、房颤/房扑和高血压等事件率在ZR组较低；其中感染4.69比7.80，嗜中性粒细胞减少0.78比9.14。暴露校正率描述单位时间事件频率，不能取代患者层面的累计风险、停药率、剂量调整、出血及心血管事件细节。ZR长期持续治疗的依从性、停药原因及高龄合并症患者获益风险，仍是正式展示和上市后监测的重点。

\subsection*{（四）摘要与ASH正式展示的时点差异}
截至2026年11月10日，ASH在线摘要已经公开上述预设主要分析结果；它并非仅公布HR的简短标题，而包含研究设计、主要疗效、安全性和OS成熟度。企业称完整研究结果将在2026年12月13日ASH全会报告。现阶段不应预写尚未公开的更新数字。后续应逐项核对数据截止日期、随访时长、PFS评审口径、OS事件数、严重不良事件分项、停药率、生活质量和关键亚组；若数据截止时间或分析集不同，应明确标注为更新分析，避免将摘要与现场展示拼接成同一数据截点。

\section*{三、竞争格局与审批路径}
\subsection*{（一）竞争重点是治疗策略和证据可比性}
一线MCL并非空白市场。美国FDA已于2025年批准阿卡替尼联合BR用于不适合ASCT的初治患者，ECHO三期研究以BR为对照，中位随访49.8个月时PFS HR为0.73；但该方案保留苯达莫司汀化疗。SHINE研究评估伊布替尼联合BR，TRIANGLE则在较年轻、适合强化治疗的患者中评估含伊布替尼的免疫化疗和移植策略。另有无化疗双联或三联的早期研究，但证据等级和患者范围不同。跨试验在入组年龄、风险分层、对照方案、随访、后续治疗及终点评估方面有差异，不能据HR、缓解率或单臂ORR作疗效排名。

\RUCTableCaption{表1 一线MCL主要证据与治疗路径核对}
\begin{center}\zihao{5}\linespread{1.0}\selectfont
\begin{tabularx}{0.98\textwidth}{>{\centering\arraybackslash}p{2.25cm} >{\centering\arraybackslash}p{2.15cm} X >{\centering\arraybackslash}p{2.3cm}}
\toprule
方案/研究 & 人群与阶段 & 关键证据 & 研判边界 \\
\midrule
ZR / MANGROVE & 年长或不适合ASCT；全球随机三期 & 510例；中位随访44.1个月；PFS HR 0.57；4年PFS 62.4\% vs 46.2\% & 无化疗；持续泽布替尼；OS未成熟；未获批一线 \\
阿卡替尼+BR / ECHO & 初治且不适合ASCT；随机三期，FDA已批 & 598例；中位随访49.8个月；PFS HR 0.73 & 仍含化疗；不同研究不可直接比较 \\
伊布替尼+BR / SHINE & 年长初治患者；随机三期 & 支持BTKi加入BR的疾病控制证据 & 化疗与BTKi暴露并存；需结合当期指南及安全性 \\
含伊布替尼强化方案 / TRIANGLE & 较年轻、适合强化治疗；随机三期 & 比较强化免疫化疗、移植及伊布替尼策略 & 目标人群与MANGROVE不同 \\
ZR无化疗组合 / II期探索 & 初治患者；单臂或小样本研究 & 可提供可行性与初步活性信号 & 不等同随机三期确证证据 \\
\bottomrule
\end{tabularx}
\end{center}
\RUCNote{注：表中各研究来自不同试验，不构成头对头疗效比较。BR为苯达莫司汀联合利妥昔单抗；ASCT为自体造血干细胞移植；HR为风险比。资料来源：ASH摘要、FDA及公开临床研究。}
\RUCTableAfterSpace

\subsection*{（二）监管进度从申报受理走向审评结论}
美国FDA于11月4日宣布受理泽布替尼联合利妥昔单抗用于成人初治MCL的补充新药申请并给予优先审评。公司同时披露已向EMA递交申请，并获中国国家药品监督管理局药品审评中心（CDE）受理；申报均基于MANGROVE。三个市场的受理状态不应被描述为上市许可。应继续跟踪FDA审评完成日期及标签范围、EMA审评意见与欧盟委员会决定、NMPA技术审评及批准公告。企业披露日本一线MCL申报预计于2026年第四季度，属计划节点而非已完成事项。

\subsection*{（三）竞争优势取决于标签和实际治疗定位}
若监管认可ZR用于所研究的初治患者，差异化将来自化疗规避、口服持续治疗，以及随机三期头对头比较BR所提供的证据。但标签最终范围、指南推荐等级、对不同风险亚组的定位，以及医生对长期口服BTKi的接受度，均影响实际渗透。既有BTKi、含BTKi免疫化疗、固定疗程组合与细胞治疗在不同疾病阶段和患者中的角色会并存；新适应症放量更可能是患者分层后的治疗选择重构，而非对所有方案的整体替代。

\section*{四、业绩与估值验证}
\subsection*{（一）销售基数与新增适应症拆分}
BeOne披露BRUKINSA 2026年第三季度全球销售额为14亿美元，同比增长34\%；美国销售额为10亿美元。公司2026全年总收入指引上调至69亿至70亿美元，但总收入指引同时包含其他产品和业务，不应作为一线MCL销售基线。公司未单独披露MCL适应症销售额，因此适应症增量需独立建模，并与现有CLL/SLL、WM、复发难治MCL等适应症区分。

\subsection*{（二）收入敏感性：用患者覆盖和持续时间拆解}
在尚无公开目标市场净价、获批标签及真实世界治疗时长的情况下，不宜给出单点销售预测。可用以下公式做滚动校验：某年度新增销售额约等于“符合标签的年度新治疗患者数 × 采用ZR比例 × 平均治疗年限 × 年均净销售额/患者”。以下仅为假设情景，每年以1,000名符合标签的新治疗患者为测算单位，反映稳定治疗队列中的年化销售规模；它不是疾病流行病学估算、公司指引或峰值销售预测。高渗透情景也没有计入患者支付能力及同类治疗竞争的约束。

\RUCTableCaption{表2 一线MCL新增销售情景敏感性（每年1,000名符合标签患者）}
\begin{center}\zihao{5}\linespread{1.0}\selectfont
\begin{tabularx}{0.95\textwidth}{>{\centering\arraybackslash}p{2cm} >{\centering\arraybackslash}p{2cm} >{\centering\arraybackslash}p{2.1cm} >{\centering\arraybackslash}p{2.2cm} X}
\toprule
情景 & ZR采用率 & 平均用药时间 & 年均净销售额/患者 & 稳态年化增量 \\
\midrule
谨慎 & 20\% & 1.5年 & 7万美元 & 2,100万美元 \\
基准 & 40\% & 2.0年 & 10万美元 & 8,000万美元 \\
积极 & 60\% & 2.5年 & 13万美元 & 1.95亿美元 \\
\bottomrule
\end{tabularx}
\end{center}
\RUCNote{注：公式为年度患者数×采用率×平均用药年限×年均净销售额/患者；表内参数为分析情景假设，非公司预测或现行价格。结果未计汇率、停药、折扣、税费、地区销售分成、患者跨期治疗队列爬坡及新适应症间替代。若按单一新患者队列首年计算，情景收入分别约为420万、2,000万和7,800万美元；稳态数值反映多个年度治疗队列重叠。}
\RUCTableAfterSpace

公司的销售兑现应按三个阶段观察：获批后首先是适用患者中使用率和启动治疗人数；随后是平均用药持续时间、剂量调整与停药；最终体现为净价、报销覆盖和各市场销售分布。由于泽布替尼一线方案采用持续治疗，较长PFS一方面可能扩大在治患者存量，另一方面也提高累计药费和长期安全性审视的重要性。价格折扣、保险准入、患者援助和竞争性方案可能压低模型中的净价或采用率。

\RUCTableCaption{表3 临床证据—监管节点—销售假设核对表}
\begin{center}\zihao{5}\linespread{1.0}\selectfont
\begin{tabularx}{0.98\textwidth}{>{\centering\arraybackslash}p{1.8cm} X X >{\centering\arraybackslash}p{2cm}}
\toprule
链条节点 & 已有事实 & 下一步核对项 & 当前状态 \\
\midrule
临床证据 & MANGROVE随机三期；PFS HR 0.57；摘要中位随访44.1个月 & ASH正式展示更新截点、OS成熟度、安全性分项、生活质量、亚组 & PFS阳性；OS未成熟 \\
美国监管 & sNDA获FDA优先审评 & 审评结论、标签适用人群、剂量和警示 & 审评中 \\
欧洲监管 & EMA已受理申报 & 审评意见、欧盟批准决定和上市时间 & 审评中 \\
中国监管 & CDE已受理申报 & 技术审评结论、批准范围、医保及院内准入 & 审评中 \\
商业采用 & 2026Q3 BRUKINSA全球销售约14亿美元 & 新治疗患者、渗透率、用药时长、净价及支付覆盖 & 尚无一线MCL独立销售额 \\
\bottomrule
\end{tabularx}
\end{center}
\RUCNote{注：监管状态截至2026年11月10日，依据企业公告；受理和优先审评均不代表批准。收入预测只应在适应症、价格及覆盖率信息逐步明确后滚动更新。}
\RUCTableAfterSpace

\clearpage
\section*{五、投资策略与风险提示}
\subsection*{（一）投资策略}
建议将该事件视为泽布替尼全球化商业兑现的边际跟踪因素，而非独立的估值锚。临床端关注ASH正式展示对摘要主要结果的补充及OS成熟进度；监管端逐市场区分受理、审评、批准和商业上市；商业端以适用患者数、使用率、平均治疗时间、净价及报销为核心变量，建立实际季度数据与情景假设的桥接。若批准后公司单独披露新适应症销售或患者数据，可更新模型；在此之前不把总销售增长归因于一线MCL。

估值讨论应以现有多适应症收入和盈利能力为底座，把潜在MCL增量按获批时间、区域、渗透与净价折现，并对审评延后、适用范围窄化和治疗时长下降做敏感性分析。监管节点与临床结果是催化剂，但实际价值取决于增量收入能否在毛利、销售费用及竞争影响后转化为利润和现金流。

\subsection*{（二）风险提示}
\noindent\textbf{审批与长期证据：}优先审评不代表获批；监管机构可能要求补充分析、限制标签或不予批准，且PFS获益未必转化为OS改善。\par
\noindent\textbf{安全性与竞争：}持续BTK抑制剂暴露的感染、出血、心律失常和高血压管理及停药风险需关注；含BTKi免疫化疗、其他靶向组合和后续新方案可能改变一线选择，不同试验不可直接比较。\par
\noindent\textbf{商业化与估值：}药价、保险覆盖、指南更新、可及患者、治疗时长及实际渗透率可能低于假设；泽布替尼总销售还受其他适应症、汇率、竞争和研发投入影响，单一适应症催化可能已部分反映在预期中。
\clearpage
\section*{\centering\songti\fontsize{22}{26.4}\selectfont\textbf{参考文献}}
\RUCReference{[1]}{BeOne Medicines. BeOne Medicines Announces Third Quarter 2026 Financial Results and Business Updates[EB/OL]. (2026-11-05)[2026-11-10]. \href{https://hkexir.beonemedicines.com/news/beone-medicines-announces-third-quarter-2026-financial-results-and-business-updates/090a271d-e1d2-4195-b14c-36ffab8c3614}{公司公告链接}.}
\RUCReference{[2]}{American Society of Hematology. 6 Superior progression-free survival with zanubrutinib plus rituximab vs bendamustine-rituximab in treatment-naïve mantle cell lymphoma: Primary analysis of a phase 3 study (MANGROVE)[EB/OL]. (2026)[2026-11-10]. \href{https://onlineprogram.hematology.org/paper-details.aspx?pid=695083}{ASH摘要链接}.}
\RUCReference{[3]}{BeOne Medicines. U.S. FDA Grants Priority Review to BeOne Medicines’ BRUKINSA Combination in First-Line Mantle Cell Lymphoma[EB/OL]. (2026-11-04)[2026-11-10]. \href{https://ir.beonemedicines.com/news/us-fda-grants-priority-review-to-beone-medicines-brukinsa-combination-in-first-line-mantle-cell-lymphoma/b074b7a8-a9b4-48ce-b5ca-4b2c63fedece}{FDA优先审评公告链接}.}
\RUCReference{[4]}{U.S. Food and Drug Administration. FDA approves acalabrutinib with bendamustine and rituximab for previously untreated mantle cell lymphoma[EB/OL]. (2025-01-16)[2026-11-10]. \href{https://www.fda.gov/drugs/cancer-approvals/fda-approves-acalabrutinib-bendamustine-and-rituximab-previously-untreated-mantle-cell-lymphoma}{FDA批准公告链接}.}
\RUCReference{[5]}{National Cancer Institute. Mantle Cell Lymphoma Treatment (PDQ®)[EB/OL]. (更新日期不详)[2026-11-10]. \href{https://www.cancer.gov/types/lymphoma/hp/mantle-cell-lymphoma-treatment}{NCI治疗综述链接}.}
\RUCReference{[6]}{U.S. National Library of Medicine. Study of Zanubrutinib Plus Rituximab Compared With Bendamustine Plus Rituximab in Previously Untreated Mantle Cell Lymphoma (NCT04002297)[EB/OL]. (登记信息)[2026-11-10]. \href{https://clinicaltrials.gov/study/NCT04002297}{试验登记链接}.}
\RUCReference{[7]}{BeOne Medicines. BeOne Medicines Announces Positive Phase 3 Results for BRUKINSA in Frontline Mantle Cell Lymphoma[EB/OL]. (2026)[2026-11-10]. \href{https://ir.beonemedicines.com/news/beone-medicines-announces-positive-phase-3-results-for-brukinsa-in-frontline-mantle-cell-lymphoma/44597eac-a44b-451f-a021-48d345cdb85e}{研究顶线公告链接}.}
\RUCReference{[8]}{American Society of Hematology. ASH Annual Meeting and Exposition 2026[EB/OL]. (2026)[2026-11-10]. \href{https://www.hematology.org/meetings/annual-meeting}{ASH年会信息链接}.}
\end{document}
